% !TeX program = lualatex
% !TeX encoding = utf8
% !TeX spellcheck = uk_UA
% !TeX root =../PyscfBook.tex

%% --------------------------------------------------------
\section{Багатореферентна конфігураційна взаємодія (MRCI)}
%% --------------------------------------------------------

\textbf{MRCI} (\textit{Multi-Reference Configuration Interaction}) --- це високоточний \textbf{пост-CASSCF} метод, призначений для врахування \textbf{динамічної електронної кореляції} в системах із \textbf{вираженим статичним (нединамічним) кореляційним ефектом}.


На відміну від \textbf{однореферентних} методів (HF, MP2, CISD, CCSD), які базуються на \emph{одному} референсному детермінанті, \texttt{MRCI} використовує \textbf{багатодетермінантний референсний простір}, сформований зазвичай за допомогою розрахунку \texttt{CASSCF}.

\subsection{Математична форма хвильової функції}

Хвильова функція \texttt{MRCI} має вигляд:
\begin{equation}
|\Psi_{\text{MRCI}}\rangle = \sum_{I \in \text{ref}} c_I |\Phi_I\rangle
+ \sum_{I \in \text{ref}} \sum_{r<a} c_{I,r}^{a} |\Phi_{I,r}^{a}\rangle
+ \sum_{I \in \text{ref}} \sum_{r<s<a<b} c_{I,rs}^{ab} |\Phi_{I,rs}^{ab}\rangle + \cdots
\end{equation}

де:
\begin{itemize}
    \item $\{|\Phi_I\rangle\}$ --- \textbf{референсні детермінанти} з \texttt{CASSCF}-простору (включають усі важливі ближньовироджені конфігурації),
    \item $|\Phi_{I,r}^{a}\rangle$, $|\Phi_{I,rs}^{ab}\rangle$ --- \textbf{одинарні (S)} та \textbf{подвійні (D)} збудження з референсного простору в зовнішній (віртуальний) простір,
    \item $c_I, c_{I,r}^{a}, \ldots$ --- коефіцієнти, що визначаються \textbf{варіаційно}.
\end{itemize}

\subsection{Фізичний зміст: статична та динамічна кореляція}

\begin{itemize}
    \item \textbf{Статична (нединамічна) кореляція} --- враховується \textbf{всередині референсного простору} (через змішування конфігурацій у \texttt{CASSCF}).
    \item \textbf{Динамічна кореляція} --- додається \textbf{через збудження з референсного простору} в зовнішні орбіталі.
\end{itemize}

Такий двоетапний підхід дозволяє MRCI точно описувати:
\begin{itemize}
    \item потенціальні енергетичні поверхні (PES),
    \item спектроскопічні константи збуджених станів,
    \item реакційні профілі з розривом зв’язків,
    \item системи з сильним багатореферентним характером (\ce{O3}, \ce{Cr2}, \ce{FeO} тощо).
\end{itemize}

\subsection{Основні варіанти MRCI}

\begin{itemize}
    \item \texttt{MRCISD} --- одинарні та подвійні збудження (найпоширеніший),
    \item \texttt{MRCISD+Q} --- додає \textbf{квадрупольну корекцію} (за Девідсоном або подібними формулами),
    \item \texttt{MRCISD(T)} --- включає потрійні збудження (надвисока точність, але $\mathcal{O}(N^{10})$).
\end{itemize}

\subsection{Реалізація в \texttt{PySCF}}

У пакеті \texttt{PySCF} \texttt{MRCI} доступний через модуль \texttt{mrci}, який працює \textbf{після} \texttt{CASSCF}-оптимізації. Приклад для \ce{H2} (розрив зв’язку):

\begin{minted}{python}
from pyscf import gto, scf, mcscf, mrci

# Молекула
mol = gto.M(atom='H 0 0 0; H 0 0 1.4', basis='cc-pvdz', spin=0)

# HF
mf = scf.RHF(mol).run()

# CASSCF: 2 електрони в 2 орбіталях (σ, σ*)
cas = mcscf.CASSCF(mf, ncas=2, nelecas=2)
cas.run()

# MRCISD
mrci_calc = mrci.MRCI(cas)
mrci_calc.run()

print(f"E(MRCISD) = {mrci_calc.e_tot:.10f} Eh")
\end{minted}

Можливості \texttt{PySCF} MRCI:
\begin{itemize}
    \item підтримка \inlinecode{nroots} --- розрахунок кількох станів одночасно,
    \item симетрія (\inlinecode{mol.symmetry = True}),
    \item спінова адаптація,
    \item експорт CI-векторів і густин.
\end{itemize}


\begin{commentbox}
\textbf{MRCI --- «золотий стандарт» точності}, але лише для \textbf{малих систем} ($\leq 12$--$15$ атомів).
Для більших молекул краще використовувати \texttt{CASPT2}.
\end{commentbox}
